\section{Introduction}
\label{sec:introduction}

Delays in escalating the care of critically ill patients in the emergency department (ED) are harmful: ED patients who waited six hours or
more for transfer to the intensive care unit (ICU) had a higher in-hospital mortality~\cite{chalfin2007delayed}. EDs rank arriving
patients by triage acuity, and early warning scores such as the National Early Warning Score 2 (NEWS2)~\cite{rcp2017news2} add a rule-based score computed from vital signs.
The MIMIC-IV-ED benchmark~\cite{xie2022mimic4ed} made machine-learning comparisons on triage data, which predict critical care outcomes better than the triage acuity alone~\cite{raita2019triage}, reproducible on public data.

Whether such a model would help depends on when its score is computed. The benchmark predicts its critical outcome at triage, but its pipeline also provides every vital sign charted during the stay, and a model that reads all of them scores the stay at ED departure. Part of that record is produced after the decision the score is meant to inform: how often a patient was measured and how long the patient
stayed reflect the course of the stay. A model can therefore reach a high area under the receiver operating characteristic curve (AUROC) for a reason that is of
no use at the bedside. Leakage of this kind is a recognized cause of optimistic results in machine-learning studies~\cite{kapoor2023leakage}. Landmarking restricts both the inputs and the
population at risk to what holds at a chosen time~\cite{vanhouwelingen2007landmarking}, and so states the prediction time explicitly.

A score that is acted on also needs a statement of how far it can be trusted. Conformal prediction attaches to each patient a set of labels that contains the true
label at a chosen rate, with a finite-sample guarantee~\cite{lei2018distributionfree,vovk2022algorithmic}. The ED setting strains this guarantee in two ways. The
guarantee assumes that calibration patients and later patients are exchangeable, which a change in case mix or practice over time can break; adaptive methods such as
Error-quantified Conformal Inference (ECI)~\cite{wu2025eci} update the calibration online for this reason. The guarantee is also marginal: when few patients deteriorate,
a set can meet its coverage over all patients while rarely containing the deterioration label for those who need it. For a department, the relevant quantities are how
many deteriorating patients a set flags and how many false alarms it raises to do so.

A recorded literature search (Section~\ref{sec:related-work} and Supplementary Table~\ref{S-tab:search}) found restricted-input benchmarks for
MIMIC-IV-ED~\cite{solaiman2026earlydeterioration}. It found no study that measured what quantities that describe how the stay was monitored add to arrival information in an end-of-stay model, and none that evaluated conformal sets for ED deterioration at a stated prediction time together with their false-alarm cost.

The aim of this study is to determine how well an ED deterioration model discriminates when its inputs are restricted to a stated prediction time, and what its
conformal sets guarantee for the patients who deteriorate. The use we evaluate is a score available to ED clinicians two hours after arrival, after triage and before
disposition. Using MIMIC-IV-ED~\cite{johnson2023mimicived} with a chronological, patient-level split, we audit the inputs of an end-of-stay model against their source timestamps, redefine the task at a landmark two hours after arrival with a 24-hour outcome window, compare a gated recurrent unit with decay (GRU-D)~\cite{che2018grud}, gradient-boosted trees~\cite{chen2016xgboost}, and logistic regression under an equal tuning budget, and calibrate the primary model with static, adaptive, class-conditional, and periodically recalibrated conformal methods.

This study makes three contributions.
\begin{itemize}
\item An input-timing audit of an end-of-stay model, a design that the benchmark pipeline makes easy. It finds retained vital-sign rows charted after ED departure and separates the discrimination of the model into arrival information, workflow quantities (ED duration and measurement counts, which contain no physiological value), and vital-sign values, so the end-of-stay task does not support a claim about early warning.
\item A comparison of recurrent, tree-based, and linear models at a two-hour landmark under a pre-specified tuning budget and five pre-specified paired contrasts: timing-aware trees discriminate slightly better than the recurrent model in this cohort and tuning budget, and the gain of the recurrent model over a model of the vital-sign sequence alone comes from the static profile.
\item An evaluation of conformal calibration for the patients who deteriorate: sets with marginal coverage rarely contain the deterioration label, class-conditional sets restore it, and the positive predictive value and false alarms of the resulting sets lie on the precision--recall curve of the score, across coverage levels, label delays, and a controlled drift, with online updating compared with periodic recalibration.
\end{itemize}
We do not propose a new conformal method. The cohort is a retrospective single-center proxy, and the results concern prediction and calibration behavior, not clinical
benefit.
